%% LyX 2.4.4 created this file.  For more info, see https://www.lyx.org/.
%% Do not edit unless you really know what you are doing.
\documentclass[english]{article}
\usepackage[T1]{fontenc}
\usepackage[latin9]{inputenc}
\usepackage{enumitem}
\usepackage{amsmath}
\usepackage{amsthm}
\usepackage{amssymb}
\usepackage{geometry}
\geometry{verbose,tmargin=1in,bmargin=1in,lmargin=1in,rmargin=1in}

\makeatletter
%%%%%%%%%%%%%%%%%%%%%%%%%%%%%% Textclass specific LaTeX commands.
\numberwithin{equation}{section}
\numberwithin{figure}{section}
\newlength{\lyxlabelwidth}      % auxiliary length 
\theoremstyle{definition}
\newtheorem*{defn*}{\protect\definitionname}
\theoremstyle{plain}
\newtheorem*{lem*}{\protect\lemmaname}
\theoremstyle{plain}
\newtheorem*{thm*}{\protect\theoremname}

%%%%%%%%%%%%%%%%%%%%%%%%%%%%%% User specified LaTeX commands.
\usepackage{fancyhdr}
\usepackage{lastpage}
\pagestyle{fancy}
\usepackage{enumitem}
\usepackage{ifthen}
\everymath=\expandafter{\the\everymath\displaystyle}
%\DeclareMathSizes{10}{10}{10}{10}
\usepackage{multicol}

\usepackage{tikz}
\usetikzlibrary{patterns}
\usetikzlibrary{plotmarks}
\usepackage{pgfplots}
\pgfplotsset{compat=newest}

\usepackage{calc}
\usepackage[nomessages]{fp}

\usepackage{datetime}
% Added by lyx2lyx
\setlength{\parskip}{\smallskipamount}
\setlength{\parindent}{0pt}

\makeatother

\usepackage{babel}
\providecommand{\definitionname}{Definition}
\providecommand{\lemmaname}{Lemma}
\providecommand{\theoremname}{Theorem}

\begin{document}
\renewcommand{\author}{Dr.\,James Ashe}

\newcommand{\classNum}{335}

\newcommand{\classSec}{A}

\newcommand{\nameType}{\hspace{4cm} Solutions}

\newcommand{\examType}{Homework 6}

\newcommand{\Date}{Due date: \formatdate{26}{4}{2013}} %%\AdvanceDate[12] \today

%\newdate{my_date}{\day}{\month}{\year}%   
%\AdvanceDate[-5] 
%\displaydate{my_date}

%%First page's header%%%%%%%%%%%%%%%%%%%%%%
\fancypagestyle{firststyle} %%header settings for first pagr
{
	\fancyhead[L]{MTH \classNum{}(\classSec{}) \author{}\\
	\textbf{\examType{}} \Date{}}
	\fancyhead[C]{\textbf{\nameType}}
	\setlength{\headsep}{14pt}
}
%%%%%%%%%%%%%%%%%%%%%%%%%%%%%%%%%%%%%%%%%%

%%Next page's header%%%%%%%%%%%%%%%%%%%%%%%
%\setlist{nolistsep} %eliminates space between list items
\lhead{MTH \classNum{}(\classSec{})} 
\chead{\textbf{\nameType}} 
\cfoot{\thepage{}~of~\pageref{LastPage}} 
\setlength{\headsep}{5pt} 
%%%%%%%%%%%%%%%%%%%%%%%%%%%%%%%%%%%%%%%%%%%

%%Various settings; see the preamble for other settings%%%%%
%%\setenumerate[0]{leftmargin=12pt,itemindent=0pt,labelwidth=0pt}
\renewcommand{\labelenumi}{\textbf{\arabic{enumi}.}}
\renewcommand{\labelenumii}{\textbf{\alph{enumii}.}}
\thispagestyle{firststyle}
%%%%%%%%%%%%%%%%%%%%%%%%%%%%%%%%%%%%%%%%%%

\small
\begin{defn*}
A \emph{group homomorphism} is a function $f:G_{1}\rightarrow G_{2}$
between two groups that preserves the group operation; i.e., $f(ab)=\underset{\mbox{operation of }G_{2}}{\underbrace{f(a)f(b)}\mbox{.}}$
\end{defn*}
\begin{itemize}
\item [ex.]The mapping $f:\mathbb{Z}\rightarrow\mathbb{Z}_{n}$ where $\mathbb{Z}$
and $\mathbb{Z}_{n}$ are groups under addition and modular addition
respectively. Since $f(a+b)=(a+b)\bmod n=a\bmod n+b\bmod n$.
\end{itemize}
\begin{defn*}
Let $G$ be a group a $H$ a subset of $G$. For any $a\in G$, the
set $aH=\left\{ ah|h\in H\right\} $ is called the \emph{left coset
of H in G containing a. }$Ha=\left\{ ha|h\in H\right\} $ is called
the \emph{right coset of H in G containing a}. 
\end{defn*}
\begin{lem*}
For $H$ a subgroup of $G$ and $a,b\in G$, the following properties
of cosets hold:

\begin{enumerate}
\item $a\in aH$

\begin{proof}
$e\in H$ since $H\leq G$. So $a=ae\in aH$.
\end{proof}
\item $aH=H$ if and only if $a\in H$.

\begin{proof}
Let $aH=H$. Then $a=ae\in aH=H$. Suppose that $a\in H$. $aH\subseteq H$
since $H$ is closed, and $H\subseteq aH$ since if $h\in H$ then
$a^{-1}h\in H$. Thus $h=eh=\left(aa^{-1}\right)h=a\left(a^{-1}h\right)\in aH$. 
\end{proof}
\item $aH=bH$ or $aH\cap bH=\emptyset$

\begin{proof}
Suppose that $aH\cap bH\neq\emptyset$. Then there exists an $x\in aH\cap bH$,
and there are $h_{1},h_{2}\in H$ such that $x=ah_{1}=bh_{2}$. Thus
$a=xh_{1}^{-1}=bh_{2}h_{1}^{-1}$. So then $aH=\left(bh_{2}h_{1}^{-1}\right)H=bH$
since $h_{2},h_{1}^{-1}\in H$ (because $H\leq G$). 
\end{proof}
\item $\left|aH\right|=\left|Ha\right|=\left|bH\right|=\left|H\right|$.

\begin{proof}
Two sets have equal order if and only if there is a bijection between
them. Let $h\mapsto ah$ define a map from $H$ to $aH$. $\left|aH\right|\leq\left|H\right|$
by definition, so the $h\mapsto ah$ is onto. Because $G$ is a group
there exists an $a^{-1}$. So then there exists an inverse map $a^{-1}(ah)\mapsto h$
so that map is $1-1$. Thus it is a bijection and $\left|H\right|=\left|aH\right|$.
Similarly a bijection exists for $bH$ and $Ha$, and $\left|aH\right|=\left|Ha\right|=\left|bH\right|=\left|H\right|$.
\end{proof}
\end{enumerate}
\end{lem*}
\begin{thm*}
\textbf{Lagrange's Theorem: $\left|H\right||\left|G\right|$. }If
$G$ is a finite group and $H$ is a subgroup of $G$, then $\left|H\right|$
divides $\left|G\right|$. (also $\frac{\left|G\right|}{\left|H\right|}$is
the number of distinct left cosets in $G$).

\begin{proof}
Let $a_{1}H,a_{2}H,\dots,a_{n}H$ be a list of distinct left cosets
of $G$. For each $a\in G$, $aH=a_{i}H$ for some $1\leq i\leq n$.
By part 1 of the lemma, $a\in aH=a_{i}H$. Thus each $a\in G$ is
in on of the cosets from the list, and we can write 
\[
G=a_{1}H\cup a_{2}H\cup\dots\cup a_{n}H\mbox{.}
\]
By part 3 of the lemma, $a_{i}H\cap a_{j}H=\emptyset$ for each $i$
and $j$ from the list. So then 
\[
\left|G\right|=\left|a_{1}H\right|+\left|a_{2}H\right|+\dots+\left|a_{n}H\right|\mbox{.}
\]
By part 4, $\left|a_{i}H\right|=\left|H\right|$. So we have that
$\left|G\right|=n\left|H\right|$ which implies that $\left|H\right||\left|G\right|$. 
\end{proof}
\end{thm*}
\begin{enumerate}
\item Prove that if $\left|G\right|\leq4$ then $G$ is abelian. Then prove
that if $\left|G\right|<6$ then $G$ is abelian. 

\begin{proof}
Suppose that $\left|G\right|\leq4$ and that $ab\neq ba$ for some
$a,b\in G$. Since $ab\neq ba$ then $a\neq b$ and neither $a$ or
$b$ are the identity. So then $e,a,b,ab,ba$ are all distinct elements
in $G$ but then $\left|G\right|>4$ which is a contradiction. 

Alternatively, since $G$ is nontrivial it has an element $a$ such
that $a\neq e$. By Lagrange's Theorem, $\left|a\right|=4\mbox{ or }2$.
If $\left|a\right|=4$ then $G$ is cyclic and thus abelian. If all
of its nontrivial elements have order 4 then $e=ab\left(ab\right)\Rightarrow ba=ba\left(abab\right)=b\left(aa\right)bab=bbab=ab$.
Since $aa=e$ and $bb=e$. 

For the second part we need only show $G$ is abelian when $\left|G\right|=5$.
$5$ is prime so then if $a\in G$ and $a\neq e$ then$\left|a\right||5$.
So then $G$ is cyclic. 
\end{proof}
\end{enumerate}
\begin{enumerate}[resume]
\item Let $f:G_{1}\rightarrow G_{2}$ be a group homomorphism and let $H\leq G_{2}$.
Prove that $f^{-1}(H)=\left\{ g\in G_{1}|f(g)\in H\right\} $ is a
subgroup of $G_{1}$.

\begin{proof}
$e\in f^{-1}(H)$ so $f^{-1}(H)\neq\emptyset$. Also $f^{-1}(H)\subseteq G_{1}$
by definition. Now let $a,b\in f^{-1}(H)$. By definition of $f^{-1}(H)$,
$f(a),f(b)\in H$, and thus $f(b^{-1})\in H$ since $H$ is a subgroup.
So then $f(ab^{-1})=f(a)f(b)^{-1}\in H$ since $f$ is a homomorphism
and $H$ is a subgroup, which implies that $ab^{-1}\in f^{-1}(H)$.
\end{proof}
\end{enumerate}

\end{document}
